\chapterbackmatter{Solutions to Weinberg Exercises}
\ExerciseEditorialNote

\begin{enumerate}
  \WeinbergSolution{1}
  Put
  \[
    v_n\equiv\langle\phi_n\rangle_{\rm VAC},
    \qquad
    f_\alpha=i\phi_n(t_\alpha)_{nm}v_m.
  \]
  Under an infinitesimal gauge transformation,
  \(\delta\phi_n=i\epsilon_\beta(t_\beta)_{n\ell}\phi_\ell\), so
  \[
    \frac{\delta f_\alpha(x)}{\delta\epsilon_\beta(y)}
      =-(t_\alpha v)_n(t_\beta\phi(x))_n\,
        \delta^4(x-y).
  \]
  The Faddeev--Popov part of the Lagrangian is therefore, up to the
  overall ghost-sign convention,
  \[
    \boxed{\mathcal L_{\rm gh}
      =-\omega_\alpha^*(t_\alpha v)_n
        (t_\beta\phi)_n\omega_\beta.}
  \]
  Writing \(\phi=v+\phi'\), the quadratic kernel is
  \[
    - (t_\alpha v)_n(t_\beta v)_n
      =\mu^2_{\alpha\beta},
  \]
  the gauge-boson mass-squared matrix of
  Eq.~\eqref{eq:21.1.7}.  Hence, on the broken-generator subspace,
  \[
    \boxed{\Delta^{\rm gh}_{\alpha\beta}(k)
      =(\mu^2)^{-1}_{\alpha\beta}}
  \]
  with no momentum dependence.  Zero modes of \(\mu^2\) signal that this
  scalar condition has not fixed the unbroken gauge transformations and
  require an additional gauge condition.

  Although the ghost density is a local polynomial, its constant
  propagator does not fall as \(k^{-2}\).  Insertions of the
  \(\phi'\omega^*\omega\) vertex therefore generate ultraviolet
  divergences of unbounded degree.  This generalized unitarity gauge is
  not power-counting renormalizable.  It can be used as a formal gauge
  for physical amplitudes, but renormalizability is made manifest only
  in a finite-\(\xi\) gauge such as that of Section~\ref{sec:21.2}.

  \WeinbergSolution{2}
  A real \(SU(2)\) triplet has weak isospin \(T=1\) and hypercharge
  \(Y=0\).  Choose
  \(\langle\vec\phi\rangle=(0,0,v_T)\).  Its covariant derivative gives
  \[
    \frac12(D_\mu\vec\phi)^2
      \supset\frac12g^2v_T^2
        \bigl[(A_\mu^1)^2+(A_\mu^2)^2\bigr],
  \]
  while \(A_\mu^3\) and the hypercharge field \(B_\mu\) do not acquire
  masses.  Thus
  \[
    m_W^2=g^2v_T^2,\qquad
    \mathcal M^2_{\rm neutral}=0.
  \]
  The breaking pattern is
  \[
    SU(2)_L\times U(1)_Y
      \longrightarrow U(1)_{T_3}\times U(1)_Y,
  \]
  not the observed single \(U(1)_{\rm em}\).  There would be two
  massless neutral gauge bosons and no massive \(Z^0\).

  There is a second obstruction.  A renormalizable charged-lepton
  Yukawa interaction must contain a scalar doublet, because
  \(\bar L e_R\) transforms as a doublet.  A real triplet cannot contract
  \(\bar L e_R\) to an \(SU(2)\) singlet, so the usual quark and charged
  lepton masses would not be generated.  A triplet can supplement a
  doublet, but it cannot replace it.  If both have expectation values,
  a \(Y=0\) triplet shifts \(m_W\) without the corresponding shift of
  \(m_Z\), producing at tree level
  \(\rho=m_W^2/(m_Z^2\cos^2\theta_W)>1\); precision data therefore force
  such a triplet expectation value to be small.

  \WeinbergSolution{3}
  Define \(a_\mu=(g_\mu-2)/2\) and
  \(s_W=\sin\theta_W\).  The \(Z\mu\bar\mu\) couplings may be written in
  terms of
  \[
    v_\mu=-1+4s_W^2,\qquad a_\mu^{(Z)}=-1,
  \]
  after extracting the common weak coupling.  Expanding the one-loop
  vertex integral through leading order in \(m_\mu^2/m_Z^2\) gives
  \[
    \boxed{a_\mu^{Z}
      =\frac{G_Fm_\mu^2}{8\sqrt2\,\pi^2}
       \left[-\frac53+\frac13(1-4s_W^2)^2\right]
      +O\!\left(\frac{m_\mu^4}{m_Z^4}
        \ln\frac{m_Z^2}{m_\mu^2}\right).}
  \]
  The large negative term is the axial coupling; the vector term is
  strongly suppressed because \(1-4s_W^2\) is small.

  The physical neutral scalar \(h\) couples as
  \(-m_\mu h\bar\mu\mu/v\), with \(v^{-2}=\sqrt2G_F\).  Its exact
  one-loop Feynman-parameter representation is
  \[
    \boxed{a_\mu^h
      =\frac{G_Fm_\mu^2}{4\sqrt2\,\pi^2}
        \int_0^1dx\,
        \frac{x^2(2-x)}
        {x^2+(m_h^2/m_\mu^2)(1-x)}.}
  \]
  For \(m_h\gg m_\mu\),
  \[
    a_\mu^h
      =\frac{G_Fm_\mu^2}{4\sqrt2\,\pi^2}
       \frac{m_\mu^2}{m_h^2}
       \left[\ln\frac{m_h^2}{m_\mu^2}-\frac76\right]
       +O\!\left(\frac{m_\mu^4}{m_h^4}
         \ln\frac{m_h^2}{m_\mu^2}\right).
  \]
  It is positive and Yukawa-suppressed.  As a check, adding the charged
  \(W\) contribution
  \(a_\mu^W=(G_Fm_\mu^2/8\sqrt2\pi^2)(10/3)\) gives, apart from the tiny
  scalar term,
  \[
    a_\mu^{\rm EW,1\,loop}
      =\frac{G_Fm_\mu^2}{8\sqrt2\,\pi^2}
       \left[\frac53+\frac13(1-4s_W^2)^2\right].
  \]

  \WeinbergSolution{4}
  The electromagnetic vertex of an on-shell charged spin-one particle
  is conventionally parameterized by \(g_1(0)\), \(\kappa(0)\), and
  \(\lambda(0)\), with magnetic dipole moment
  \[
    \boldsymbol\mu
      =\frac{e}{2m}\,[g_1(0)+\kappa(0)+\lambda(0)]\,\mathbf S.
  \]
  The Yang--Mills kinetic term fixes at tree level
  \[
    g_1^\gamma=1,\qquad \kappa_\gamma=1,\qquad\lambda_\gamma=0.
  \]
  Therefore the \(W^+\) has gyromagnetic ratio \(g_W=2\) and
  \[
    \boxed{\mu_{W^+}=\frac{e}{m_W}}
  \]
  along a unit spin projection.  The sign reverses for \(W^-\).

  The \(Z^0\) is electrically neutral, and the renormalizable
  electroweak Yang--Mills Lagrangian contains no diagonal
  \(\gamma Z Z\) vertex.  Its lowest-order diagonal magnetic dipole
  moment is consequently
  \[
    \boxed{\mu_{Z^0}=0.}
  \]
  This does not exclude loop-induced transition form factors involving
  neutral gauge bosons, but those are not tree-level magnetic moments.

  \WeinbergSolution{5}
  A complete sequential generation makes the same \(n_g\)-dependent
  contribution to the three inverse-coupling equations once hypercharge
  is given its unified normalization.  This is visible directly in
  Eqs.~\eqref{eq:21.5.9}--\eqref{eq:21.5.11}: the \(n_g\) terms cancel
  both in
  \[
    \frac1{g_s^2}-\frac1{g^2}
  \]
  and in
  \[
    \frac1{g^2}-\frac3{5g'^2}.
  \]
  Hence the one-loop predictions
  \eqref{eq:21.5.15} for \(\sin^2\theta_W\) and
  \eqref{eq:21.5.16} for the unification scale \(M\) are independent of
  \(n_g\).  The discovery of a fourth complete generation would
  therefore leave those two predictions unchanged at this accuracy.

  It would not be entirely invisible.  Solving any one of the individual
  running equations shows that the extra matter lowers
  \(\alpha_U^{-1}\), so the unified coupling is larger.  Gauge-boson
  mediated proton decay, whose rate scales schematically as
  \(\alpha_U^2/M^4\), would then be enhanced if \(M\) were fixed.
  Threshold corrections from the new fermion masses and higher-loop
  running would also modify precision predictions.  The cancellation
  relies on adding a complete anomaly-free generation; incomplete
  multiplets or exotic quantum numbers would generally change the
  relative running and hence the predicted meeting of the couplings.

  \WeinbergSolution{6}
  Let condensates of charges \(q_i\) have nonzero magnitudes.  Far from
  boundaries, minimizing the gradient energy gives
  \[
    \boldsymbol\nabla\theta_i-\frac{q_i}{\hbar}\mathbf A=0.
  \]
  Around a closed contour, single-valuedness requires
  \[
    \Delta\theta_i=2\pi n_i,\qquad
    \Phi=\oint\mathbf A\cdot d\mathbf x
      =\frac{2\pi\hbar\,n_i}{q_i}
      \quad\hbox{for every }i.
  \]
  Thus allowed fluxes form the intersection of the lattices
  \((2\pi\hbar/q_i)\mathbb Z\).  For commensurate charges this
  intersection is another nonzero lattice.  If some ratio \(q_i/q_j\)
  is irrational, the only common value is
  \[
    \boxed{\Phi=0.}
  \]
  Nonzero topological flux sectors, the usual universal flux quantum,
  and isolated stable vortices would therefore disappear.  Persistent
  currents could relax by continuously redistributing the relative
  phases rather than only by jumps of one common flux quantum.

  Local effects that depend only on a mass for the photon, such as the
  Meissner effect and zero dc resistance, can remain: the London mass is
  proportional to \(\sum_iq_i^2|v_i|^2\).  Junction currents, however,
  can contain harmonics with Josephson frequencies
  \[
    \nu_i=\frac{|q_i\Delta V|}{2\pi\hbar}.
  \]
  Incommensurate charges give no single fundamental period; the current
  is generally quasiperiodic, and independent relative phases supply
  additional collective modes unless interactions lock them.

  Finally, exact incommensurate charges are incompatible with the
  compact electromagnetic \(U(1)\) assumed in Section~\ref{sec:21.6}:
  no common nonzero gauge parameter acts trivially on every field.  The
  hypothesis therefore changes the global gauge group as well as the
  phenomenology, which is precisely why ordinary charge
  commensurability underlies flux quantization.
\end{enumerate}
